% 概要
\section{概要}

\subsection{プロジェクト名}
Embedded Module Architecture — ESP32-S3 CAM サンプルプロジェクト

\subsection{コンセプト}
本プロジェクトは、組み込みシステムにおける\textbf{Moduleインターフェース設計パターン}のリファレンス実装である。
ESP32-S3 N16R8 CAM開発ボードを題材に、カメラ・センサー・LED・LCD等の多様なハードウェアモジュールを
統一的なインターフェースで管理する手法を示す。

\subsection{設計思想}

\begin{funcblock}{基本理念：どんなプロジェクトにも流用できる汎用的なコアクラス}
IModule・ModuleTimerは\textbf{プロジェクトに一切依存しない}汎用コンポーネントとして設計する。
C++テンプレートを活用し、プロジェクト固有のデータ型は利用側で自由に指定できる。
コアクラスをコピーするだけで、どんな組み込みプロジェクトにもそのまま流用可能である。
\end{funcblock}

本アーキテクチャの根幹となる3つの要素：

\begin{enumerate}
    \item \textbf{IModule<T>テンプレートインターフェース} — C++テンプレートにより、プロジェクト固有の型に依存せず、
    全モジュールが\texttt{init()}と\texttt{update(T\& data)}を実装する統一的なライフサイクル管理
    \item \textbf{ModuleTimerクラス} — \texttt{delay()}を一切使用しないノンブロッキングなタイマー制御
    \item \textbf{SystemDataパターン} — プロジェクト固有の共有状態をサブ構造体で一元管理し、モジュール間のデータ連携を実現
\end{enumerate}

この設計により、以下を実現する：
\begin{itemize}
    \item \textbf{コアクラスの完全な汎用性}（プロジェクト非依存）
    \item モジュールの追加・削除が容易
    \item 処理順序が決定論的（ループ内で順番に実行）
    \item 複数回出力や処理の重複が発生しない
    \item ゲームループ（ECS）に類似した、組み込みに適した実行モデル
    \item 参照渡し（\texttt{\&}）によるnullポインタ問題の排除
\end{itemize}
